\documentclass[11pt,a4paper]{article}
\usepackage[utf8]{inputenc}
\usepackage[T1]{fontenc}
\usepackage{lmodern}
\usepackage[spanish,es-noshorthands,es-nodecimaldot,es-tabla]{babel}
\usepackage{amsmath,amssymb}
\usepackage{booktabs}
\usepackage{array}
\usepackage[margin=2.3cm]{geometry}
\usepackage[hidelinks]{hyperref}

\newcommand{\J}{\mathcal{J}}
\newcommand{\Sg}{\mathcal{S}}
\newcommand{\N}{\mathcal{N}}
\newcommand{\M}{\mathcal{M}}
\DeclareMathOperator*{\lexmin}{lex\,min}

\title{Modelo matemático\\[2pt]
\large Flow shop híbrido con sublotes consistentes para la producción de guantes industriales}
\date{}

\begin{document}
\maketitle

\section{Descripción y supuestos}

Un conjunto de órdenes de trabajo debe pasar, en el mismo orden, por cuatro
etapas (troquelado, ensamble, costura y volteado), cada una con máquinas
paralelas no relacionadas. Cada orden se divide en un número fijo de sublotes
de tamaño variable que conservan su tamaño en todas las etapas
(\emph{sublotes consistentes}). Se busca minimizar la tardanza total y, como
criterio secundario, el makespan.

\begin{enumerate}
  \item Todas las máquinas están disponibles desde $t=0$. El tiempo se mide en
        horas hábiles del calendario laboral (sin noches, domingos ni festivos).
  \item Cada máquina procesa un sublote a la vez y no hay interrupciones.
  \item Cada sublote se procesa en una sola máquina por etapa. Sublotes de la
        misma orden pueden procesarse a la vez en máquinas distintas.
  \item Un sublote puede iniciar la etapa $s+1$ en cuanto termina la etapa $s$
        (transferencia inmediata, sin tiempo de transporte y con almacenamiento
        intermedio ilimitado).
  \item Todas las máquinas de una etapa pueden procesar todas las referencias,
        con tiempos distintos.
  \item En troquelado hay un alistamiento cada vez que cambia la referencia en
        una máquina y antes del primer sublote de cada máquina. El alistamiento
        no depende del sublote y puede hacerse mientras la máquina espera
        su llegada.
  \item El tiempo de proceso es proporcional al tamaño del sublote, y los
        tamaños son números enteros de docenas.
\end{enumerate}

\section{Conjuntos e índices}

\begin{center}
\begin{tabular}{>{$}l<{$} p{12cm}}
\toprule
j, j' \in \J & Órdenes de trabajo. \\
s \in \Sg = \{1,\dots,|\Sg|\} & Etapas, en el orden de la ruta ($s=1$: troquelado). \\
m \in \M_s & Máquinas de la etapa $s$. \\
b \in \{1,\dots,B_j\} & Sublotes de la orden $j$. \\
i, i' \in \N & Conjunto de todos los sublotes; $i=(j,b)$, $i'=(j',b')$. \\
0 & Sublote ficticio que representa el inicio de la secuencia de una máquina. \\
\bottomrule
\end{tabular}
\end{center}

\section{Parámetros}

\begin{center}
\begin{tabular}{>{$}l<{$} p{12.5cm}}
\toprule
Q_j & Cantidad de la orden $j$ [docenas]. \\
\rho_j & Referencia (código de producto) de la orden $j$. \\
r_j & Liberación de la orden $j$ [h hábiles]: horas acumuladas al final del día anterior a la fecha de liberación. \\
d_j & Fecha de entrega de la orden $j$ [h hábiles]: horas acumuladas al final del día de entrega. \\
a_{jm} & Tiempo de proceso por docena de la orden $j$ en la máquina $m$ [h/docena]:
         $a_{jm} = \tau_{\rho_j m} / 60$, con $\tau_{\rho m}$ el tiempo estándar [min/docena]. \\
q^{\min} & Tamaño mínimo de sublote [docenas] ($q^{\min}=1$). \\
B_j & Número de sublotes de la orden $j$:
      $B_j = \max\bigl\{1,\ \min\{\bar B,\ \lfloor Q_j/q^{\min} \rfloor\}\bigr\}$, con $\bar B = 10$. \\
\sigma_{ii'} & Alistamiento en troquelado si $i$ precede inmediatamente a $i'$ [h]:
               $\sigma_{ii'} = \sigma$ si $\rho_j \neq \rho_{j'}$ y $0$ si no;
               $\sigma_{0i'} = \sigma$ (montaje inicial). $\sigma = 4{,}98/60$ h.
               En las etapas $s\ge 2$ no hay alistamiento. \\
H & Constante grande (cota del horizonte), p.\,ej.
    $H = \max_j r_j + \sum_{s}\sum_{j} Q_j \max_{m\in\M_s} a_{jm} + |\N|\,\sigma$. \\
\bottomrule
\end{tabular}
\end{center}

\section{Variables de decisión}

\begin{center}
\begin{tabular}{>{$}l<{$} p{12.5cm}}
\toprule
x_{jb} \in \mathbb{Z}_{\ge 0} & Tamaño del sublote $b$ de la orden $j$ [docenas], igual en todas las etapas. \\
y_{ism} \in \{0,1\} & 1 si el sublote $i$ se procesa en la máquina $m$ de la etapa $s$. \\
z_{ii'sm} \in \{0,1\} & 1 si el sublote $i$ precede inmediatamente al sublote $i'$ en la máquina $m$ de la etapa $s$
                        ($i=0$: $i'$ es el primero de la máquina). \\
S_{is},\ C_{is} \ge 0 & Inicio y terminación del sublote $i$ en la etapa $s$ [h]. \\
T_j \ge 0 & Tardanza de la orden $j$. \\
C_{\max} \ge 0 & Makespan. \\
\bottomrule
\end{tabular}
\end{center}

\section{Función objetivo}

Objetivo lexicográfico: primero la tardanza total y, entre soluciones con la
misma tardanza, el menor makespan:
\begin{equation}
  \lexmin\ \Bigl(\ \sum_{j\in\J} T_j\ ,\ C_{\max}\Bigr)
  \label{eq:obj}
\end{equation}
En un solucionador se puede usar el equivalente
$\min \sum_j T_j + \varepsilon\, C_{\max}$ con $\varepsilon>0$ suficientemente pequeño.

\section{Restricciones}

\paragraph{Tamaño de los sublotes (sublotes consistentes).}
\begin{align}
  \sum_{b=1}^{B_j} x_{jb} &= Q_j
    && \forall j\in\J \label{eq:suma}\\
  x_{jb} &\ge q^{\min}
    && \forall j\in\J,\ b=1,\dots,B_j \label{eq:min}
\end{align}
Como $x_{jb}$ no lleva índice de etapa, el sublote tiene el mismo tamaño en
todas las etapas.

\paragraph{Asignación de máquinas.}
\begin{align}
  \sum_{m\in\M_s} y_{ism} &= 1
    && \forall i\in\N,\ s\in\Sg \label{eq:asig}
\end{align}

\paragraph{Tiempo de proceso.}
\begin{align}
  C_{is} &\ge S_{is} + a_{jm}\, x_{jb} - H\,(1 - y_{ism})
    && \forall i=(j,b)\in\N,\ s\in\Sg,\ m\in\M_s \label{eq:fin}
\end{align}
Si el sublote se asigna a la máquina $m$, termina después de su tiempo de
proceso en esa máquina; para las demás máquinas la restricción queda
desactivada. En el óptimo se cumple con igualdad, porque un $C_{is}$ mayor solo
retrasa el resto de la programación.

\paragraph{Liberación y precedencia entre etapas.}
\begin{align}
  S_{i1} &\ge r_j
    && \forall i=(j,b)\in\N \label{eq:rel}\\
  S_{i,s+1} &\ge C_{is}
    && \forall i\in\N,\ s=1,\dots,|\Sg|-1 \label{eq:prec}
\end{align}
La restricción \eqref{eq:prec} se aplica a cada sublote por separado: el
sublote $b$ pasa a la etapa siguiente sin esperar a los demás sublotes de su
orden.

\paragraph{Secuencia en cada máquina.}
Cada sublote asignado a una máquina tiene exactamente un antecesor inmediato
(otro sublote o el inicio ficticio $0$) y a lo sumo un sucesor; cada máquina
tiene a lo sumo un primer sublote:
\begin{align}
  \sum_{\substack{i\in\{0\}\cup\N\\ i\neq i'}} z_{ii'sm} &= y_{i'sm}
    && \forall s\in\Sg,\ m\in\M_s,\ i'\in\N \label{eq:ant}\\
  \sum_{\substack{i'\in\N\\ i'\neq i}} z_{ii'sm} &\le y_{ism}
    && \forall s\in\Sg,\ m\in\M_s,\ i\in\N \label{eq:suc}\\
  \sum_{i'\in\N} z_{0i'sm} &\le 1
    && \forall s\in\Sg,\ m\in\M_s \label{eq:primero}
\end{align}

\paragraph{No solapamiento y alistamientos.}
\begin{align}
  S_{i'1} &\ge C_{i1} + \sigma_{ii'} - H\,(1 - z_{ii'1m})
    && \forall m\in\M_1,\ i,i'\in\N,\ i\neq i' \label{eq:disy1}\\
  S_{i'1} &\ge \sigma_{0i'}\, z_{0i'1m}
    && \forall m\in\M_1,\ i'\in\N \label{eq:setup0}\\
  S_{i's} &\ge C_{is} - H\,(1 - z_{ii'sm})
    && \forall s\ge 2,\ m\in\M_s,\ i,i'\in\N,\ i\neq i' \label{eq:disy}
\end{align}
Si $i$ precede inmediatamente a $i'$, el sublote $i'$ empieza después de que
termina $i$, más el alistamiento por cambio de referencia en troquelado. El
alistamiento se cuenta desde que la máquina queda libre, por lo que se
superpone con la espera del sublote: el inicio efectivo es
$\max\{C_{i1} + \sigma_{ii'},\ r_{j'}\}$, que resulta de combinar
\eqref{eq:disy1} con \eqref{eq:rel}. Como todo sublote tiene tamaño positivo
(por \eqref{eq:min}), \eqref{eq:disy1} y \eqref{eq:disy} impiden ciclos en la
secuencia.

\paragraph{Tardanza y makespan.}
La orden $j$ termina cuando su último sublote sale de la última etapa,
$C_j = \max_b C_{(j,b),|\Sg|}$; por eso basta con:
\begin{align}
  T_j &\ge C_{(j,b),|\Sg|} - d_j
    && \forall j\in\J,\ b=1,\dots,B_j \label{eq:tj}\\
  C_{\max} &\ge C_{i,|\Sg|}
    && \forall i\in\N \label{eq:cmax}
\end{align}

\paragraph{Dominio de las variables.}
\begin{align}
  & x_{jb} \in \mathbb{Z}_{\ge 0};\quad
    y_{ism},\, z_{ii'sm} \in \{0,1\};\quad
    S_{is},\, C_{is},\, T_j,\, C_{\max} \ge 0 \label{eq:dom}
\end{align}

\paragraph{Conversión a fechas.}
Si $\Phi(\delta)$ son las horas hábiles acumuladas al final del día $\delta$,
la fecha en que termina la orden $j$ es
$\min\{\delta : \Phi(\delta) \ge C_j\}$.

\section{Resolución con el algoritmo genético}

Con 10 órdenes de 10 sublotes ($|\N| = 100$) y 36 máquinas, el modelo tiene del
orden de $\sum_s |\M_s|\,|\N|^2 \approx 3{,}6\times 10^5$ variables binarias
$z_{ii'sm}$, por lo que se resuelve con un algoritmo genético cuyo
decodificador construye soluciones factibles de
\eqref{eq:suma}--\eqref{eq:dom}.

\paragraph{Cromosoma.} $\Omega = (\pi, \Lambda)$, donde
\begin{itemize}
  \item Nivel 1: $\pi_j = (\pi_{j1},\dots,\pi_{jB_j})$, con $\pi_{jb} > 0$ y
        $\sum_b \pi_{jb} = 1$: proporciones de los tamaños de los sublotes de la orden $j$.
  \item Nivel 2: $\Lambda_s$, una permutación de $\N$ para cada etapa $s$:
        prioridad de los sublotes en la etapa.
\end{itemize}

\paragraph{Decodificación del nivel 1 (tamaños).}
Con $L_j = Q_j - B_j\, q^{\min}$ docenas libres,
\[
  x_{jb} = q^{\min} + \lfloor \pi_{jb} L_j \rfloor + \delta_{jb},
\]
donde $\delta_{jb}\in\{0,1\}$ reparte las
$L_j - \sum_b \lfloor \pi_{jb} L_j \rfloor$ docenas restantes entre los sublotes con
mayor parte fraccionaria (método del mayor residuo). Así se cumplen
\eqref{eq:suma}--\eqref{eq:min}.

\paragraph{Decodificación del nivel 2 (programación).}
Para cada etapa $s = 1,\dots,|\Sg|$ y cada sublote $i=(j,b)$ en el orden de $\Lambda_s$:
\begin{enumerate}
  \item Instante de llegada: $e_{is} = r_j$ si $s=1$, o $e_{is} = C_{i,s-1}$ si $s\ge 2$.
  \item Para cada $m \in \M_s$, con $p_{ism} = a_{jm}\, x_{jb}$:
    \begin{itemize}
      \item Etapa 1 (con alistamiento): si $F_m$ es el fin del último sublote
            programado en $m$ y $\rho^{\text{últ}}_m$ su referencia,
            $S_{ism} = \max\{F_m + \sigma\,\mathbb{1}[\rho^{\text{últ}}_m \neq \rho_j],\ e_{is}\}$
            (para la primera asignación de $m$: $F_m = 0$ y se cobra $\sigma$).
      \item Etapas $s\ge 2$: $S_{ism}$ es el primer instante $\ge e_{is}$ en
            que $m$ tiene un hueco libre de duración $p_{ism}$.
    \end{itemize}
  \item Se elige la máquina con menor terminación (regla ECT):
        $m^* = \arg\min_{m\in\M_s} \{S_{ism} + p_{ism}\}$, y se fija
        $C_{is} = S_{ism^*} + p_{ism^*}$.
\end{enumerate}
Luego $T_j = \max\{0,\ \max_b C_{(j,b),|\Sg|} - d_j\}$ y la aptitud
es el par $\bigl(\sum_j T_j,\ C_{\max}\bigr)$, comparado lexicográficamente.

\paragraph{Operadores.}
\begin{itemize}
  \item Población inicial aleatoria: $\pi_j$ con distribución de Dirichlet$(1,\dots,1)$
        y $\Lambda_s$ permutaciones aleatorias. Un individuo se reemplaza por
        la mejor regla de despacho (FIFO, EDD, SPT, LPT, MST, CR) con sublotes iguales.
  \item Selección por torneo de tamaño 3 y elitismo de un individuo.
  \item Cruce (probabilidad $0{,}90$):
    \begin{itemize}
      \item Nivel 1 (aritmético): para cada orden, con $\lambda \sim U(0,1)$,
            $\pi_j^{h_1} = \lambda\pi_j^{p_1} + (1-\lambda)\pi_j^{p_2}$ y
            $\pi_j^{h_2} = (1-\lambda)\pi_j^{p_1} + \lambda\pi_j^{p_2}$.
      \item Nivel 2 (cruce de orden, OX): en cada etapa, el hijo conserva un
            tramo aleatorio de $\Lambda_s^{p_1}$ en las mismas posiciones y
            completa el resto en el orden de $\Lambda_s^{p_2}$ (y al revés para el otro hijo).
    \end{itemize}
  \item Mutación (probabilidad $0{,}50$ por hijo):
    \begin{itemize}
      \item Nivel 2: con probabilidad $0{,}30$, intercambio de dos sublotes en $\Lambda_s$ de una etapa al azar.
      \item Nivel 1: para cada orden, con probabilidad $0{,}40$,
            $\pi_{jb} \leftarrow \max\{10^{-9},\ \pi_{jb} + \mathcal{N}(0,\,0{,}15^2)\}$
            para un sublote $b$ al azar, seguido de normalización.
    \end{itemize}
  \item Parada: 50 generaciones sin mejora o 1000 generaciones. Población: 150.
\end{itemize}

\section{Instancia (BD\_HFS.xlsx)}

\begin{center}
\begin{tabular}{lr}
\toprule
Órdenes $|\J|$ & 10 \\
Etapas y máquinas & Troquelado (4), Ensamble (15), Costura (9), Volteado (8) \\
Sublotes por orden $B_j$ & 10 \\
Tamaño de sublote & docenas enteras, $q^{\min} = 1$ docena \\
Alistamiento en troquelado $\sigma$ & 4,98 min \\
Liberaciones $r_j$ & 0 a 119 h hábiles \\
Fechas de entrega $d_j$ & 400 a 672 h hábiles \\
\midrule
Mejor regla de despacho (EDD) & $\sum T_j = 145{,}00$ h, \ $C_{\max} = 817{,}00$ h \\
Algoritmo genético (semilla 42) & $\sum T_j = 0$ h, \ $C_{\max} = 647{,}33$ h \\
\bottomrule
\end{tabular}
\end{center}

\section*{Referencia}
Gürsoy Yılmaz, B., \& Yılmaz, Ö. F. (2022). Lot streaming in hybrid flowshop
scheduling problem by considering equal and consistent sublots under machine
capability and limited waiting time constraint. \emph{Computers \& Industrial
Engineering}, 173, 108745.

\end{document}
